% ============================================================
%  CareerPath AI — Document de Proiectare Arhitecturală
%  Template academic conform cerințelor universitare
%  Compile: pdflatex proiectare_arhitecturala.tex (x2)
% ============================================================
\documentclass[12pt, a4paper]{report}

% ── Encoding & Language ─────────────────────────────────────
\usepackage[T1]{fontenc}
\usepackage[utf8]{inputenc}
\usepackage[romanian]{babel}

% ── Font ────────────────────────────────────────────────────
\usepackage{lmodern}
\usepackage{microtype}

% ── Page Layout ─────────────────────────────────────────────
\usepackage[
  top=2.5cm, bottom=2.5cm,
  left=3cm,  right=2.5cm,
  headheight=15pt
]{geometry}
\usepackage{fancyhdr}
\usepackage{emptypage}
\usepackage{setspace}
\onehalfspacing

% ── Colors ──────────────────────────────────────────────────
\usepackage[dvipsnames, table]{xcolor}
\definecolor{primary}{HTML}{2C3E8C}
\definecolor{accent}{HTML}{1A6EBD}
\definecolor{lightblue}{HTML}{EBF3FB}
\definecolor{darkgray}{HTML}{2D2D2D}
\definecolor{rowgray}{HTML}{F5F5F5}
\definecolor{success}{HTML}{1E7E34}
\definecolor{muted}{HTML}{6C757D}

% ── Graphics ────────────────────────────────────────────────
\usepackage{tikz}
\usetikzlibrary{
  shapes.geometric, shapes.rounded corners,
  arrows.meta, positioning, fit, calc,
  backgrounds, shadows, decorations.pathreplacing,
  matrix
}
\usepackage{pgfplots}
\pgfplotsset{compat=1.18}

% ── Tables ──────────────────────────────────────────────────
\usepackage{booktabs}
\usepackage{tabularx}
\usepackage{multirow}
\usepackage{makecell}
\usepackage{array}
\renewcommand{\arraystretch}{1.3}

% ── Boxes ───────────────────────────────────────────────────
\usepackage{tcolorbox}
\tcbuselibrary{skins, breakable}

% ── Lists ───────────────────────────────────────────────────
\usepackage{enumitem}

% ── Hyperlinks ───────────────────────────────────────────────
\usepackage[hidelinks, unicode]{hyperref}
\hypersetup{
  pdftitle={CareerPath AI - Document de Proiectare Arhitecturala},
  pdfauthor={Rares Constantinescu, Francis, Vlad},
  bookmarksnumbered=true
}

% ── Section styling ─────────────────────────────────────────
\usepackage{titlesec}
\titleformat{\chapter}[block]
  {\normalfont\large\bfseries\color{primary}\MakeUppercase}
  {\thechapter.}{0.6em}{}[{\titlerule[0.8pt]}]
\titlespacing{\chapter}{0pt}{0pt}{12pt}

\titleformat{\section}
  {\normalfont\normalsize\bfseries\color{primary}}
  {\thesection.}{0.5em}{}
\titlespacing{\section}{0pt}{14pt}{6pt}

\titleformat{\subsection}
  {\normalfont\normalsize\bfseries}
  {\thesubsection.}{0.5em}{}
\titlespacing{\subsection}{0pt}{10pt}{4pt}

% ── Header/Footer ───────────────────────────────────────────
\pagestyle{fancy}
\fancyhf{}
\fancyhead[L]{\small\color{muted} Document de Proiectare Arhitecturală — CareerPath AI}
\fancyhead[R]{\small\color{muted} \leftmark}
\fancyfoot[C]{\small\thepage}
\renewcommand{\headrulewidth}{0.4pt}

% ── Captions ────────────────────────────────────────────────
\usepackage{caption}
\captionsetup{
  font=small, labelfont=bf,
  labelsep=period, justification=centering
}

% ── Code ────────────────────────────────────────────────────
\usepackage{listings}
\lstset{
  basicstyle=\ttfamily\footnotesize,
  keywordstyle=\color{accent}\bfseries,
  commentstyle=\color{muted}\itshape,
  stringstyle=\color{success},
  breaklines=true,
  frame=tb,
  rulecolor=\color{primary!30},
  backgroundcolor=\color{lightblue!40},
  numbers=left, numberstyle=\tiny\color{muted},
  xleftmargin=15pt, xrightmargin=5pt,
  aboveskip=8pt, belowskip=8pt
}

% ── Misc ────────────────────────────────────────────────────
\usepackage{graphicx}
\usepackage{float}
\usepackage{pifont}
\usepackage{fontawesome5}

% ── Custom note box ─────────────────────────────────────────
\newtcolorbox{notabox}[1][]{
  colback=lightblue!60, colframe=primary!60,
  arc=4pt, boxrule=0.6pt,
  left=6pt, right=6pt, top=4pt, bottom=4pt,
  fonttitle=\small\bfseries\color{primary},
  #1
}

% ============================================================
\begin{document}
\selectlanguage{romanian}

% ── COPERTĂ ─────────────────────────────────────────────────
\begin{titlepage}
\begin{tikzpicture}[remember picture, overlay]
  \fill[primary] (current page.north west)
    rectangle ([yshift=-3.5cm]current page.north east);
  \fill[accent!70!primary] ([yshift=-3.5cm]current page.north west)
    rectangle ([yshift=-3.8cm]current page.north east);
\end{tikzpicture}

\vspace*{1.8cm}
\begin{center}
  {\color{white}\Large\bfseries UNIVERSITATEA — FACULTATEA DE INFORMATICĂ}\\[0.3cm]
  {\color{white!80!gray}\normalsize Proiect Software — 2025/2026}
\end{center}

\vspace{3cm}

\begin{center}
  {\color{primary}\rule{14cm}{1.5pt}}\\[0.8cm]
  {\Huge\bfseries\color{darkgray} CareerPath AI}\\[0.4cm]
  {\Large\color{accent} Platformă Inteligentă de Dezvoltare a Carierei}\\[0.6cm]
  {\color{primary}\rule{14cm}{1.5pt}}\\[0.8cm]
  {\LARGE\bfseries\color{darkgray} DOCUMENT DE PROIECTARE ARHITECTURALĂ}
\end{center}

\vspace{2.5cm}

\begin{center}
\begin{tabular}{ll}
  \textbf{Versiune:}   & 1.0 \\[4pt]
  \textbf{Data:}       & Aprilie 2026 \\[4pt]
  \textbf{Autori:}     & Rareș Constantinescu, Francis, Vlad \\[4pt]
  \textbf{Stare:}      & În dezvoltare \\[4pt]
  \textbf{Clasificare:}& Uz intern \\
\end{tabular}
\end{center}

\vfill
\begin{center}
  {\small\color{muted} Acest document descrie arhitectura componentelor implementate\\
  până în Aprilie 2026 în cadrul proiectului CareerPath AI.}
\end{center}
\end{titlepage}

% ── CUPRINS ─────────────────────────────────────────────────
\tableofcontents
\clearpage

% ============================================================
\chapter{Introducere}
% ============================================================

\section{Scopul sistemului}

CareerPath AI este o platformă web inteligentă orientată pe evenimente, destinată
sprijinirii utilizatorilor în procesul de dezvoltare a carierei profesionale. Sistemul
utilizează modele de limbaj de mari dimensiuni (LLM), căutare semantică prin vectori și
procesare asincronă bazată pe mesagerie pentru a automatiza analiza CV-urilor și
direcționarea utilizatorilor către oportunități de angajare sau trasee de învățare adaptate.

Scopul prezentului document este de a descrie arhitectura software a componentelor
implementate până la data curentă, incluzând structura subsistemelor, fluxul de date,
gestionarea persistenței și controlul accesului.

\textbf{Componente implementate la data redactării:}
\begin{itemize}[leftmargin=*]
  \item Infrastructura Docker (PostgreSQL, Redis, Qdrant, Redpanda, Keycloak)
  \item Frontend — interfață web (Next.js 14): autentificare, înregistrare, dashboard
  \item API Gateway — punct unic de intrare (NestJS + Fastify)
  \item User Service — gestionarea utilizatorilor (Java Spring Boot + Keycloak)
  \item CV Analyzer Agent — analiză CV cu Gemini AI + vector embeddings (FastAPI)
\end{itemize}

\section{Definiții și Acronime}

\begin{center}
\rowcolors{2}{rowgray}{white}
\begin{tabularx}{\textwidth}{lX}
\toprule
\rowcolor{primary!15}
\textbf{Termen / Acronim} & \textbf{Definiție} \\
\midrule
LLM & Large Language Model — model de limbaj de mari dimensiuni \\
RAG & Retrieval-Augmented Generation — generare augmentată prin recuperare \\
JWT & JSON Web Token — token compact pentru autentificare \\
RS256 & Algoritm de semnare JWT bazat pe RSA cu SHA-256 \\
API & Application Programming Interface — interfață de programare a aplicației \\
REST & Representational State Transfer — stil arhitectural pentru servicii web \\
ORM & Object-Relational Mapping — maparea obiectelor pe tabele relaționale \\
SSE & Server-Sent Events — protocol push unidirecțional HTTP \\
CV & Curriculum Vitae \\
Embedding & Reprezentare vectorială densă a unui text în spațiu n-dimensional \\
Qdrant & Bază de date vectorială open-source pentru căutare semantică \\
Redpanda & Platformă de streaming compatibilă Kafka \\
NestJS & Framework Node.js pentru aplicații server-side scalabile \\
FastAPI & Framework Python modern pentru construirea de API-uri \\
MinIO & Stocare obiect compatibilă S3 \\
Keycloak & Soluție open-source pentru Identity \& Access Management (OAuth2/OIDC) \\
LangGraph & Framework pentru agenți AI stateful multi-pas \\
\bottomrule
\end{tabularx}
\end{center}

\section{Documente de referință}

\begin{enumerate}[leftmargin=*]
  \item \textbf{Documentul de Specificare a Cerințelor} — CareerPath AI SRS v1.0, 2026
  \item \textbf{Documentația oficială Next.js 14} — \url{https://nextjs.org/docs}
  \item \textbf{Documentația NestJS} — \url{https://docs.nestjs.com}
  \item \textbf{Documentația FastAPI} — \url{https://fastapi.tiangolo.com}
  \item \textbf{Google Gemini API} — model \texttt{gemini-2.5-flash}, embedding \texttt{gemini-embedding-001}
  \item \textbf{Qdrant Documentation} — \url{https://qdrant.tech/documentation}
  \item \textbf{Redpanda Documentation} — \url{https://docs.redpanda.com}
  \item \textbf{Keycloak Documentation} — \url{https://www.keycloak.org/documentation}
  \item \textbf{RFC 7519} — JSON Web Token (JWT) Standard
  \item \textbf{Fișierul \texttt{docker-compose.yml}} din repository-ul proiectului — definiția infrastructurii
\end{enumerate}

% ============================================================
\chapter{Obiective de Proiectare}
% ============================================================

Arhitectura CareerPath AI a fost concepută pornind de la următoarele obiective principale
de proiectare:

\begin{enumerate}[leftmargin=*, label=\textbf{OD-\arabic*.}]

  \item \textbf{Separarea responsabilităților (Separation of Concerns)}\\
  Fiecare serviciu are o singură responsabilitate bine definită. Frontend-ul nu
  conține logică de business; API Gateway-ul nu accesează baze de date direct;
  agenții AI nu interacționează niciodată direct cu utilizatorii.

  \item \textbf{Comunicare asincronă prin evenimente}\\
  Componentele din stratul AI comunică exclusiv prin evenimente Redpanda (Kafka).
  Niciun apel HTTP direct nu există între agenți, eliminând cuplarea strânsă și
  permițând scalarea independentă a fiecărui agent.

  \item \textbf{Scalabilitate orizontală}\\
  Fiecare serviciu rulează în container Docker independent și poate fi scalat
  orizontal fără modificări arhitecturale. Redpanda suportă nativ partițiile,
  iar Qdrant suportă clustering.

  \item \textbf{Securitate prin design}\\
  Autentificarea se realizează cu chei asimetrice RS256. API Gateway-ul injectează
  identitatea utilizatorului (\texttt{x-user-id}) în toate cererile downstream —
  serviciile interne nu re-verifică niciodată tokenul JWT.

  \item \textbf{Izolarea datelor pe domenii}\\
  Fiecare serviciu accesează exclusiv tabelele sale din baza de date. Schimbul de
  date între servicii se face exclusiv prin evenimente Redpanda, nu prin acces
  direct la baza de date a altui serviciu.

  \item \textbf{Observabilitate și mentenanță}\\
  Fiecare serviciu expune un endpoint \texttt{/health} standardizat. Hot-reload
  este configurat pentru toate serviciile în modul de dezvoltare, minimizând
  ciclurile de rebuild.

  \item \textbf{Reproductibilitate a mediului}\\
  Întreaga infrastructură este definită declarativ în \texttt{docker-compose.yml}.
  Profilurile Docker (\texttt{frontend}, \texttt{backend}, \texttt{agents}) permit
  fiecărui dezvoltator să pornească exclusiv subsistemul relevant.

\end{enumerate}

% ============================================================
\chapter{Arhitectura Propusă}
% ============================================================

\section{Prezentarea generală a arhitecturii sistemului}

Sistemul CareerPath AI adoptă o arhitectură pe două straturi, bazată pe microservicii
și comunicare orientată pe evenimente:

\begin{notabox}[title={Principiu arhitectural fundamental}]
\textbf{Stratul 1} (conversațional) interacționează cu utilizatorul prin HTTP/WebSocket.
\textbf{Stratul 2} (worker agents) procesează în background și comunică exclusiv
prin topicuri Redpanda. Cele două straturi nu se apelează direct niciodată.
\end{notabox}

\vspace{0.5cm}

\textbf{Subsistemele implementate și funcționalitățile alocate:}

\begin{center}
\rowcolors{2}{rowgray}{white}
\begin{tabularx}{\textwidth}{llX}
\toprule
\rowcolor{primary!15}
\textbf{Subsistem} & \textbf{Strat} & \textbf{Funcționalități alocate} \\
\midrule
Frontend (Next.js) & Prezentare & Autentificare, înregistrare, dashboard utilizator, gestionare sesiune \\
API Gateway (NestJS) & Gateway & Rutare, autentificare JWT, CORS, proxy spre servicii interne \\
User Service (Spring) & Core & Gestionare profil utilizator, integrare Keycloak OAuth2 \\
CV Analyzer Agent & AI/Worker & Scorare CV, analiză gap, generare embeddings, publicare rezultate \\
PostgreSQL & Date & Persistență utilizatori, upload-uri CV, sesiuni \\
Redis & Date & Cache sesiuni JWT, rate limiting \\
Qdrant & Date & Stocare și căutare vectori (embeddings CV/joburi) \\
Redpanda & Mesagerie & Bus de evenimente între servicii \\
Keycloak & Securitate & Identity Provider OAuth2/OIDC, emitere tokeni \\
\bottomrule
\end{tabularx}
\end{center}

\subsection{Diagrama generală a arhitecturii}

\begin{figure}[H]
\centering
\begin{tikzpicture}[
  svc/.style={
    rounded corners=5pt, minimum width=3.4cm, minimum height=1.1cm,
    inner sep=7pt, font=\small\bfseries, align=center, drop shadow
  },
  db/.style={
    cylinder, shape border rotate=90,
    minimum width=1.6cm, minimum height=1.4cm,
    inner sep=3pt, font=\footnotesize\bfseries, align=center
  },
  arr/.style={-Stealth, thick, color=primary},
  karr/.style={-Stealth, thick, dashed, color=orange!70!black},
  lbl/.style={font=\tiny, color=muted}
]

% Stratul prezentare
\node[svc, fill=darkgray!90, text=white] (fe) at (0, 5) {Frontend\\{\footnotesize Next.js 14}};
\node[above=0.1cm of fe, font=\footnotesize\bfseries\color{muted}] {Strat Prezentare};

% Gateway
\node[svc, fill=red!15, draw=red!50] (gw) at (0, 3) {API Gateway\\{\footnotesize NestJS + Fastify}};
\node[above right=0.05cm and 0.3cm of gw, font=\footnotesize\color{muted}] {:8000};

% Core services
\node[svc, fill=blue!12, draw=blue!50] (us) at (-3.5, 1) {User Service\\{\footnotesize Spring Boot}};
\node[above right=0.05cm and 0.1cm of us, font=\footnotesize\color{muted}] {:8001};

% AI Agent
\node[svc, fill=accent!15, draw=accent] (cva) at (3.5, 1) {CV Analyzer\\{\footnotesize FastAPI + Gemini}};
\node[above right=0.05cm and 0.1cm of cva, font=\footnotesize\color{muted}] {:8006};

% Redpanda
\node[svc, fill=orange!15, draw=orange!60, minimum width=8cm] (rp) at (0, -1.2) {
  Redpanda — Event Bus (Kafka-compatible)\\{\footnotesize topicuri: cv.uploaded, cv.analyzed}
};

% Databases
\node[db, fill=blue!10, draw=blue!40] (pg) at (-5.5, -0.2) {Postgres\\:5432};
\node[db, fill=red!10, draw=red!40] (redis) at (-4, -0.2) {Redis\\:6379};
\node[db, fill=purple!10, draw=purple!40] (qd) at (5.5, -0.2) {Qdrant\\:6333};
\node[db, fill=orange!10, draw=orange!40] (kc) at (-5.5, 3) {Keycloak\\:8484};

% Arrows
\draw[arr] (fe) -- node[lbl,right]{HTTP/HTTPS} (gw);
\draw[arr] (gw) -- node[lbl,above]{\texttt{x-user-id}} (us);
\draw[arr] (gw) -- node[lbl,above]{proxy} (cva);
\draw[arr] (us) -- (pg);
\draw[arr] (us) -- (redis);
\draw[arr] (gw) to[bend right=15] node[lbl,right]{validate} (kc);
\draw[karr] (cva) -- node[lbl,right]{consume/produce} (rp);
\draw[arr] (cva) -- node[lbl,right]{store} (qd);

% Legenda
\node[arr, draw=primary, fill=primary, minimum width=0.5cm, minimum height=0.1cm] at (-1,-2.8) {};
\node[right] at (-0.7,-2.8) {\footnotesize HTTP / REST};
\node[draw=orange!70!black, fill=orange!10, minimum width=0.5cm, minimum height=0.1cm, dashed] at (2,-2.8) {};
\node[right] at (2.3,-2.8) {\footnotesize Kafka event};

\end{tikzpicture}
\caption{Diagrama generală a arhitecturii CareerPath AI (componente implementate)}
\end{figure}

\section{Descompunerea în subsisteme și responsabilități}

\subsection{Frontend — Next.js 14}

\textbf{Responsabilitate:} Interfața grafică cu utilizatorul. Gestionează sesiunea de
autentificare local prin Zustand cu persistență în \texttt{localStorage}.

\textbf{Componente implementate:}
\begin{itemize}[leftmargin=*]
  \item \texttt{/login} — formular de autentificare cu validare Zod + React Hook Form
  \item \texttt{/register} — formular de înregistrare cu validare și redirect post-succes
  \item \texttt{/dashboard} — panou principal cu guard de autentificare
  \item \texttt{auth.store.ts} — store Zustand cu persistență JWT în localStorage
  \item \texttt{api-client.ts} — client Axios configurat cu \texttt{baseURL = http://localhost:8000/api}
\end{itemize}

\textbf{Interfețe expuse:}
\begin{itemize}[leftmargin=*]
  \item Consumă: \texttt{POST /api/auth/login}, \texttt{POST /api/auth/register}, \texttt{GET /api/auth/me}
  \item Nu expune niciun endpoint propriu
\end{itemize}

\textbf{Pachete principale:}
\begin{center}
\begin{tabular}{ll}
\toprule
\textbf{Pachet} & \textbf{Rol} \\
\midrule
\texttt{next} 14 & Framework React cu App Router și SSR \\
\texttt{zustand} & State management global cu middleware persist \\
\texttt{axios} & Client HTTP cu interceptori \\
\texttt{react-hook-form} + \texttt{zod} & Validare formulare \\
\texttt{tailwindcss} + \texttt{shadcn/ui} & Sistem de design și componente UI \\
\bottomrule
\end{tabular}
\end{center}

\subsection{API Gateway — NestJS + Fastify}

\textbf{Responsabilitate:} Punct unic de intrare pentru toate cererile HTTP din frontend.
Validează JWT-ul RS256, injectează \texttt{x-user-id} și rutează cererile spre serviciile interne.

\textbf{Componente implementate:}
\begin{itemize}[leftmargin=*]
  \item \texttt{AuthController} — \texttt{POST /api/auth/login}, \texttt{POST /api/auth/register},
        \texttt{GET /api/auth/me}; semnează tokeni RS256 cu cheia privată din \texttt{/app/keys/private.pem}
  \item \texttt{UsersController} — proxy transparent spre \texttt{cp-user-service:8001}
  \item \texttt{JwtAuthGuard} — verifică semnătura RS256 cu cheia publică din \texttt{/app/keys/public.pem}
  \item \texttt{app.enableCors(\{origin: http://localhost:3000\})} — CORS configurat
\end{itemize}

\textbf{Interfețe expuse:}

\begin{center}
\rowcolors{2}{rowgray}{white}
\begin{tabular}{llll}
\toprule
\rowcolor{primary!15}
\textbf{Metodă} & \textbf{Cale} & \textbf{Auth} & \textbf{Descriere} \\
\midrule
POST & \texttt{/api/auth/register} & Nu & Creează cont, returnează JWT \\
POST & \texttt{/api/auth/login} & Nu & Autentificare, returnează JWT \\
GET  & \texttt{/api/auth/me} & Da & Returnează profilul din token \\
ALL  & \texttt{/api/users/*} & Da & Proxy spre User Service \\
GET  & \texttt{/health} & Nu & Stare serviciu \\
\bottomrule
\end{tabular}
\end{center}

\subsection{User Service — Java Spring Boot}

\textbf{Responsabilitate:} Gestionarea profilurilor utilizatorilor și validarea identității
prin Keycloak. Accesează exclusiv tabelele \texttt{users} și \texttt{cv\_uploads} din PostgreSQL.

\textbf{Componente implementate:}
\begin{itemize}[leftmargin=*]
  \item Securitate OAuth2 Resource Server cu Keycloak ca Identity Provider
  \item Endpoint-uri profil utilizator (\texttt{/api/users/profile})
  \item Conexiune PostgreSQL prin Spring Data JPA
  \item Variabile de mediu pentru configurare dinamică (fără hardcodare \texttt{localhost})
\end{itemize}

\subsection{CV Analyzer Agent — FastAPI}

\textbf{Responsabilitate:} Unicul agent AI implementat. Consumă evenimentul \texttt{cv.uploaded}
din Redpanda, aplică scorarea CV prin Gemini 2.5 Flash, generează embeddings și publică
rezultatul pe \texttt{cv.analyzed}.

\textbf{Componente implementate:}
\begin{itemize}[leftmargin=*]
  \item \texttt{cv\_upload\_consumer.py} — consumator aiokafka topic \texttt{cv.uploaded}
  \item \texttt{cv\_analyzer.py} — agent Gemini cu prompt din \texttt{prompts/cv\_scoring.txt}
  \item \texttt{embedding\_tool.py} — generare embeddings cu \texttt{gemini-embedding-001} (dim=3072)
  \item \texttt{qdrant\_client.py} — stocare și căutare vectori în colecția \texttt{cv\_embeddings}
  \item \texttt{cv\_analyzed\_publisher.py} — producător aiokafka topic \texttt{cv.analyzed}
  \item \texttt{cv\_repository.py} — persistență scor în PostgreSQL
  \item \texttt{/health} — verificare stare Gemini API, Kafka producer, Qdrant, DB pool
\end{itemize}

\subsection{Distribuția subsistemelor pe platforme}

\begin{figure}[H]
\centering
\begin{tikzpicture}[
  host/.style={rounded corners=6pt, fill=gray!10, draw=gray!50,
               inner sep=12pt, minimum width=5cm},
  cont/.style={rounded corners=4pt, inner sep=6pt, font=\footnotesize\bfseries,
               minimum width=4cm, minimum height=0.8cm, align=center}
]
  % Host machine
  \node[host, minimum width=16cm, minimum height=9cm] (host) at (7,0) {};
  \node[above=0.1cm, font=\small\bfseries\color{muted}] at (host.north) {Host Machine (Windows 11 / Linux)};

  % Docker network
  \node[rounded corners=4pt, fill=blue!5, draw=blue!30,
        inner sep=10pt, minimum width=14cm, minimum height=7cm] (docker) at (7,-0.5) {};
  \node[above=0.1cm, font=\footnotesize\color{blue!60}] at (docker.north) {Docker Network: careerpath-ai\_internal-net};

  % Containers
  \node[cont, fill=darkgray!85, text=white] at (2, 1.5) {frontend:3000};
  \node[cont, fill=red!15, draw=red!40] at (6.5, 1.5) {api-gateway:8000};
  \node[cont, fill=blue!12, draw=blue!40] at (11, 1.5) {user-service:8001};
  \node[cont, fill=accent!15, draw=accent] at (2, 0) {cv-analyzer:8006};
  \node[cont, fill=blue!10, draw=blue!30] at (6.5, 0) {postgres:5432};
  \node[cont, fill=red!10, draw=red!30] at (11, 0) {redis:6379};
  \node[cont, fill=purple!10, draw=purple!40] at (2, -1.5) {qdrant:6333};
  \node[cont, fill=orange!10, draw=orange!40] at (6.5, -1.5) {redpanda:9092};
  \node[cont, fill=orange!8, draw=orange!30] at (11, -1.5) {keycloak:8484};

\end{tikzpicture}
\caption{Distribuția containerelor Docker pe platforma hardware}
\end{figure}

\section{Managementul datelor persistente}

\subsection{Sistemul de baze de date}

Platforma utilizează \textbf{patru sisteme de stocare} distincte, fiecare optimizat
pentru tipul de date gestionat:

\begin{center}
\rowcolors{2}{rowgray}{white}
\begin{tabularx}{\textwidth}{llXl}
\toprule
\rowcolor{primary!15}
\textbf{Sistem} & \textbf{Tip} & \textbf{Date gestionate} & \textbf{Port} \\
\midrule
PostgreSQL 16 & Relațional & Utilizatori, upload-uri CV, sesiuni & 5432 \\
Redis 7 & Key-Value & Cache JWT, rate limiting (TTL) & 6379 \\
Qdrant v1.10 & Vectorial & Embeddings CV (dim=3072, Cosine) & 6333 \\
Redpanda & Streaming & Evenimente între servicii (log durabil) & 9092 \\
\bottomrule
\end{tabularx}
\end{center}

\subsection{Schema conceptuală a bazei de date PostgreSQL}

\begin{figure}[H]
\centering
\begin{tikzpicture}[
  tbl/.style={rounded corners=4pt, draw=primary!60, fill=lightblue!40,
              inner sep=0, font=\footnotesize},
  hdr/.style={fill=primary!20, font=\footnotesize\bfseries, inner sep=5pt,
              minimum width=5cm},
  row/.style={inner sep=4pt, minimum width=5cm, font=\footnotesize\ttfamily},
  pk/.style={row, font=\footnotesize\ttfamily\bfseries\color{primary}},
  fk/.style={row, font=\footnotesize\ttfamily\color{accent}}
]

% Table: users
\matrix[tbl, matrix of nodes, nodes in empty cells,
  row 1/.style={hdr}, row 2/.style={pk}, row 3/.style={row},
  row 4/.style={row}, row 5/.style={row}, row 6/.style={row}
] (users) at (0,0) {
  users \\
  id : UUID PK \\
  email : VARCHAR UNIQUE \\
  name : VARCHAR \\
  role : VARCHAR \\
  created\_at : TIMESTAMP \\
};

% Table: cv_uploads
\matrix[tbl, matrix of nodes, nodes in empty cells,
  row 1/.style={hdr}, row 2/.style={pk}, row 3/.style={fk},
  row 4/.style={row}, row 5/.style={row}, row 6/.style={row}
] (cvuploads) at (7,0) {
  cv\_uploads \\
  id : UUID PK \\
  user\_id : UUID FK \\
  file\_path : TEXT \\
  score : INTEGER \\
  uploaded\_at : TIMESTAMP \\
};

% Relationship
\draw[-Stealth, thick, color=accent]
  (users.east) -- node[above, font=\tiny\color{muted}]{1 \ \ \ \ \ N} (cvuploads.west);

\end{tikzpicture}
\caption{Schema conceptuală a bazei de date PostgreSQL (tabele implementate)}
\end{figure}

\subsection{Colecțiile Qdrant}

Qdrant stochează reprezentările vectoriale ale CV-urilor pentru căutare semantică:

\begin{center}
\begin{tabular}{llll}
\toprule
\textbf{Colecție} & \textbf{Dimensiune} & \textbf{Metrică} & \textbf{Model embedding} \\
\midrule
\texttt{cv\_embeddings} & 3072 & Cosine & \texttt{gemini-embedding-001} \\
\bottomrule
\end{tabular}
\end{center}

\subsection{Strategia de persistență Redis}

Redis este utilizat exclusiv ca cache cu TTL (Time-To-Live) pentru:
\begin{itemize}[leftmargin=*]
  \item Starea sesiunilor JWT (invalidare forțată la logout)
  \item Rate limiting per IP/utilizator la nivelul API Gateway
\end{itemize}
Datele din Redis sunt volatile — pierderea lor nu afectează integritatea sistemului.

\section{Controlul accesului utilizatorilor la sistem}

\subsection{Modelul de autentificare}

Sistemul utilizează autentificare bazată pe \textbf{JWT RS256} cu chei asimetrice:

\begin{figure}[H]
\centering
\begin{tikzpicture}[
  svc/.style={rounded corners=5pt, inner sep=8pt, minimum width=3cm,
              minimum height=1cm, font=\small\bfseries, align=center},
  arr/.style={-Stealth, thick, color=primary},
  lbl/.style={font=\tiny, color=muted, midway}
]
  \node[svc, fill=gray!15, draw=gray!50] (usr) at (0,0) {\faUser\\Utilizator};
  \node[svc, fill=red!10, draw=red!40] (gw)  at (4.5,0) {API Gateway\\NestJS};
  \node[svc, fill=orange!10, draw=orange!40] (kc) at (9,0) {Keycloak\\:8484};
  \node[svc, fill=blue!10, draw=blue!40] (ds)  at (6.5,-2.5) {Serviciu\\intern};

  \draw[arr] (usr) -- node[above,lbl]{1. credentials} (gw);
  \draw[arr] (gw)  -- node[above,lbl]{2. validate} (kc);
  \draw[arr, dashed] (kc) -- node[below,lbl]{3. JWKS public key} (gw);
  \draw[arr, dashed] (gw) to[bend right=30]
    node[left,lbl]{4. JWT RS256} (usr);
  \draw[arr] (gw)  -- node[right,lbl]{5. x-user-id header} (ds);
\end{tikzpicture}
\caption{Fluxul de autentificare RS256}
\end{figure}

\textbf{Principii de securitate implementate:}
\begin{enumerate}[leftmargin=*]
  \item Cheia privată (\texttt{private.pem}) semnează tokenul — se află exclusiv în API Gateway
  \item Cheia publică (\texttt{public.pem}) verifică tokenul — distribuită tuturor serviciilor
  \item Serviciile interne \textbf{nu re-verifică} tokenul JWT — acceptă \texttt{x-user-id} injectat
  \item Tokenul expiră în \textbf{15 minute}; reîmprospătarea nu este implementată încă
  \item CORS restricționat la \texttt{http://localhost:3000} în development
\end{enumerate}

\subsection{Niveluri de acces}

\begin{center}
\rowcolors{2}{rowgray}{white}
\begin{tabular}{lll}
\toprule
\rowcolor{primary!15}
\textbf{Resursă} & \textbf{Acces anonim} & \textbf{Acces autentificat} \\
\midrule
\texttt{POST /api/auth/register} & Da & Da \\
\texttt{POST /api/auth/login}    & Da & Da \\
\texttt{GET /health}             & Da & Da \\
\texttt{GET /api/auth/me}        & Nu & Da \\
\texttt{/api/users/*}            & Nu & Da \\
\texttt{/dashboard} (frontend)   & Redirect /login & Da \\
\bottomrule
\end{tabular}
\end{center}

\section{Fluxul global al controlului}

\subsection{Fluxul de autentificare și acces dashboard}

\begin{figure}[H]
\centering
\begin{tikzpicture}[
  step/.style={rounded corners=4pt, fill=lightblue!60, draw=primary!40,
               inner sep=8pt, minimum width=4cm, minimum height=0.9cm,
               font=\small, align=center},
  dec/.style={diamond, fill=orange!15, draw=orange!60,
              inner sep=3pt, font=\small\bfseries, aspect=2.5},
  arr/.style={-Stealth, thick, color=primary},
  lbl/.style={font=\tiny\color{muted}}
]
  \node[step] (s1) at (0,0) {Utilizator accesează\\aplicația (\texttt{localhost:3000})};
  \node[dec]  (d1) at (0,-1.8) {JWT în\\localStorage?};
  \node[step] (s2) at (-3.5,-3.5) {Redirect\\la \texttt{/login}};
  \node[step] (s3) at (3.5,-3.5) {Redirect\\la \texttt{/dashboard}};
  \node[step] (s4) at (-3.5,-5.5) {\texttt{POST /api/auth/login}\\sau \texttt{/register}};
  \node[step] (s5) at (-3.5,-7.5) {API Gateway semnează\\JWT RS256};
  \node[step] (s6) at (3.5,-7.5) {Frontend apelează\\\texttt{GET /api/auth/me}};
  \node[step] (s7) at (3.5,-9.5) {Dashboard randat\\cu date utilizator};

  \draw[arr] (s1) -- (d1);
  \draw[arr] (d1) -- node[lbl,left]{Nu} (s2);
  \draw[arr] (d1) -- node[lbl,right]{Da} (s3);
  \draw[arr] (s2) -- (s4);
  \draw[arr] (s4) -- (s5);
  \draw[arr] (s5.south) -- ++(0,-0.5) -| node[lbl,above]{JWT + user} (s3.south);
  \draw[arr] (s3) -- (s6);
  \draw[arr] (s6) -- (s7);
\end{tikzpicture}
\caption{Fluxul de control autentificare și acces dashboard}
\end{figure}

\subsection{Fluxul de analiză CV (implementat parțial)}

Fluxul de analiză CV este implementat la nivel de agent, dar nu este declanșat
din frontend deoarece \texttt{cv-service} (extragere PDF) nu este încă finalizat.

\begin{figure}[H]
\centering
\begin{tikzpicture}[
  step/.style={rounded corners=4pt, inner sep=8pt, minimum width=3.8cm,
               minimum height=0.9cm, font=\small, align=center},
  arr/.style={-Stealth, thick},
  karr/.style={-Stealth, thick, dashed, color=orange!70!black}
]
  \node[step, fill=gray!15, draw=gray!50] (cv) at (0,0) {CV PDF primit\\(cv-service)};
  \node[step, fill=orange!10, draw=orange!50] (t1) at (5,0)
    {\texttt{cv.uploaded}\\topic Redpanda};
  \node[step, fill=accent!12, draw=accent] (agent) at (10,0)
    {CV Analyzer Agent\\Gemini 2.5 Flash};
  \node[step, fill=purple!10, draw=purple!40] (qd) at (10,-2)
    {Qdrant\\cv\_embeddings};
  \node[step, fill=blue!10, draw=blue!40] (pg) at (10,-4)
    {PostgreSQL\\cv\_uploads.score};
  \node[step, fill=orange!10, draw=orange!50] (t2) at (5,-2)
    {\texttt{cv.analyzed}\\topic Redpanda};

  \draw[-Stealth, thick, color=gray!60] (cv) --
    node[above,font=\tiny\color{muted}]{produce} (t1);
  \draw[karr] (t1) --
    node[above,font=\tiny\color{muted}]{consume} (agent);
  \draw[-Stealth, thick, color=purple!60] (agent) --
    node[right,font=\tiny\color{muted}]{store embedding} (qd);
  \draw[-Stealth, thick, color=blue!60] (agent) --
    node[right,font=\tiny\color{muted}]{save score} (pg);
  \draw[karr] (agent) --
    node[above,font=\tiny\color{muted}]{produce} (t2);

  % Not implemented label
  \node[fill=red!10, draw=red!40, rounded corners=3pt,
        font=\tiny\bfseries\color{red!60}, inner sep=3pt]
    at (0,0.9) {neimplementat încă};
\end{tikzpicture}
\caption{Fluxul de analiză CV prin pipeline-ul Redpanda}
\end{figure}

\section{Condițiile limită}

\subsection{Condiții de inițializare}

\begin{itemize}[leftmargin=*]
  \item \textbf{Pornire la rece:} Serviciile cu dependențe de infrastructură
        (user-service, cv-analyzer-agent) implementează retry logic — dacă PostgreSQL
        sau Redpanda nu sunt disponibile la startup, serviciul reîncearcă cu backoff
        exponențial.
  \item \textbf{Chei JWT lipsă:} Dacă fișierele \texttt{private.pem} / \texttt{public.pem}
        lipsesc, API Gateway-ul pornește dar orice request protejat returnează \texttt{500}.
        Generarea cheilor se face cu \texttt{make keys}.
  \item \textbf{Volum Redpanda corupt:} Dacă Redpanda eșuează cu exit code 132
        (instrucțiune invalidă), volumul \texttt{careerpath-ai\_redpanda\_data} trebuie
        șters și recreat.
\end{itemize}

\subsection{Condiții de oprire}

\begin{itemize}[leftmargin=*]
  \item \textbf{Oprire gracefully:} \texttt{docker compose down} trimite \texttt{SIGTERM}
        tuturor containerelor. Consumatorii Kafka (aiokafka) confirmă ultimul offset
        înainte de oprire.
  \item \textbf{Pierdere Redis:} Sesiunile active sunt pierdute; utilizatorii trebuie
        să se re-autentifice. JWT-ul rămas valid în browser funcționează până la expirare.
  \item \textbf{Pierdere Qdrant:} Embeddings-urile existente sunt pierdute; CV-urile
        trebuie re-procesate. Datele relaționale din PostgreSQL nu sunt afectate.
\end{itemize}

\subsection{Condiții de eroare}

\begin{center}
\rowcolors{2}{rowgray}{white}
\begin{tabularx}{\textwidth}{lXl}
\toprule
\rowcolor{primary!15}
\textbf{Condiție} & \textbf{Comportament} & \textbf{HTTP} \\
\midrule
JWT expirat & JwtAuthGuard returnează Unauthorized & 401 \\
JWT invalid/lipsă & JwtAuthGuard returnează Unauthorized & 401 \\
User Service indisponibil & API Gateway returnează Gateway Error & 500 \\
Gemini API key lipsă & CV Analyzer returnează eroare internă & 500 \\
Port ocupat (ex. 9000) & Container MinIO eșuează; restul neafectat & — \\
CORS origine neautorizată & Browser blochează cererea & — \\
\bottomrule
\end{tabularx}
\end{center}

\subsection{Limitări cunoscute ale implementării curente}

\begin{enumerate}[leftmargin=*, label=\textbf{L\arabic*.}]
  \item JWT-ul nu are mecanism de refresh — expiră în 15 minute, necesitând re-login
  \item \texttt{POST /api/auth/register} nu persistă utilizatorul în baza de date
        (implementarea actuală a API Gateway-ului generează token în memorie; persistența
        este responsabilitatea User Service-ului, integrare în curs)
  \item CV Analyzer Agent nu poate fi declanșat din UI deoarece cv-service (upload PDF)
        nu este finalizat
  \item MinIO (stocare fișiere CV) intră în conflict cu portul 9000 pe mașina de
        dezvoltare dacă IntelliJ IDEA este activ
\end{enumerate}

% ── Final ────────────────────────────────────────────────────
\clearpage
\begin{center}
\vspace*{3cm}
{\color{primary}\rule{10cm}{1pt}}\\[0.8cm]
{\large\bfseries\color{darkgray} CareerPath AI}\\[0.3cm]
{\color{muted} Document de Proiectare Arhitecturală — v1.0}\\[0.3cm]
{\small\color{muted} Rareș Constantinescu $\cdot$ Francis $\cdot$ Vlad $\cdot$ Aprilie 2026}\\[0.6cm]
{\color{primary}\rule{10cm}{1pt}}
\end{center}

\end{document}
